\documentclass[11pt]{article}
\usepackage[T1]{fontenc}
\usepackage{amsmath,amssymb}
\usepackage{booktabs}
\usepackage{hyperref}
\usepackage[margin=1in]{geometry}
\usepackage{caption}

\newcommand{\arm}[1]{\textsc{#1}}

\title{Glyphs Don't Compute:\\
\large Structured-Prompt Benefits Compress into a Steerable Direction under Fine-Tuning}
\author{Nam Nguyen}
\date{August 2026}

\begin{document}
\maketitle

\begin{abstract}
Earlier experiments suggested that decorating reasoning traces with rare symbolic ``phase markers'' (four alchemical glyphs marking guideline, plan, step, and conclusion) improved mathematical reasoning in Qwen2.5-7B-Instruct.
We tested this claim with a controlled six-arm study: matched fine-tuning data differing only in marker format, frozen train/validation/test splits, strict answer scoring, and preregistered contrasts.
The claim fails.
Glyph-trained models score \emph{below} models trained on the same reasoning without markers (29.4\% vs.\ 34.6\% held-out accuracy) and below models trained with meaningless dots (31.8\%).
Prompting the \emph{untrained} model with structure genuinely helps ($+28$ to $+38$ points over an unstructured prompt), but dots outperform glyphs there too: the benefit is generic structure elicitation, not anything glyph-specific.
Mechanistic follow-up explains why fine-tuning destroyed even that benefit: training compresses the structured-reasoning behavior into a single late-layer activation direction keyed to the glyphs' rare byte tokens.
The direction can be injected directly---with no glyphs in the prompt---recovering and even exceeding natural glyph performance ($+14$ points over the matched prompt), while masking the glyph positions after prompt reading costs nothing.
The symbols were never the mechanism; they were a lossy handle on a steerable direction.
\end{abstract}

\section{Background: the original observation}

An earlier ablation found that prompting the untrained Qwen2.5-7B-Instruct with glyph-structured prompts scored 76.9\% on GSM8K versus 52.4\% for a natural-language prompt.
Two problems later emerged with that measurement.
First, its answer extractor simply took the \emph{last number} appearing anywhere in the output, and its logs show 80\% of natural-prompt outputs violating the expected structure---so free-form answers were frequently scored wrong for formatting reasons alone (raw completions were not saved, so re-scoring is impossible).
Second, it never compared glyphs against equally arbitrary alternatives such as plain dots.
The present study was designed to answer, under a leakage-checked and strictly scored protocol: \textbf{do phase markers do more than decorate a chain of thought?}

\section{Experimental design}

Six LoRA adapters were trained on byte-matched renderings of the same reasoning traces, differing only in output format (Table~\ref{tab:arms}).
All arms share the base model (Qwen2.5-7B-Instruct), data order, and hyperparameters (rank~16, $\alpha=32$, one epoch, effective batch~16, lr $2\times10^{-4}$ cosine, BF16, max 2{,}048 tokens).
Checkpoints were selected on validation accuracy only; the test set (7{,}319 problems from GSM8K, SVAMP, and MATH) was frozen before any training.
Evaluation is greedy decoding with a 1{,}024-token budget and a strict parser requiring an exact final answer.

\begin{table}[h]\centering\small
\begin{tabular}{ll}
\toprule
Arm & Output format \\
\midrule
\arm{semantic} & the reasoning content, no markers \\
\arm{glyph}    & same content + fixed glyphs at phase boundaries \\
\arm{dot}      & same content + neutral dot delimiters \\
\arm{random}   & glyph identities shuffled per example \\
\arm{direct}   & final answer only \\
\arm{filler}   & meaningless filler tokens, then the answer \\
\bottomrule
\end{tabular}
\caption{The six matched training arms. Each isolates one explanation for a marker benefit.}
\label{tab:arms}
\end{table}

\section{Behavioral results}

\subsection{Trained arms on the held-out test set}

Table~\ref{tab:main} shows strict accuracy for each training arm under each prompt format.
On the matched-format diagonal---the preregistered comparison---\arm{glyph} (29.4\%) loses to \arm{semantic} (34.6\%) and to \arm{dot} (31.8\%).
\arm{filler} (17.0\%) scored below \arm{direct} (21.3\%) on validation, ruling out hidden computation from extra token positions.
Markers therefore provide no benefit over identical reasoning without them, and lose even to meaningless delimiters.

\begin{table}[h]\centering\small
\begin{tabular}{lcccc}
\toprule
Training arm & neutral & glyph & dot & headings \\
\midrule
\arm{semantic} & \textbf{.346} & .346 & .178 & .217 \\
\arm{glyph}    & .314 & \textbf{.294} & .181 & .156 \\
\arm{dot}      & \textbf{.408} & .280 & \textbf{.318} & .272 \\
\arm{random}   & \textbf{.260} & .275 & .248 & .205 \\
\bottomrule
\end{tabular}
\caption{Held-out strict accuracy, training arm $\times$ prompt format ($n=7{,}319$ per cell). Bold marks each arm's matched or best condition.}
\label{tab:main}
\end{table}

\subsection{The untrained model: structure helps, glyphs are incidental}

Evaluating the base model with no adapter under the same protocol (Table~\ref{tab:base}) shows a large, genuine benefit of structured prompting---but plain dots beat the glyphs.
The near-zero neutral score largely reflects failure to produce the required answer format, mirroring (in reverse) the artifact in the original measurement.
Critically, \emph{no trained arm beats the untrained model with a structured prompt}: the best trained cell (40.8\%) is below base-with-dots (43.4\%), and glyph training sits ten points below simply glyph-\emph{prompting} the base model.
Fine-tuning added nothing and, for glyphs, destroyed value.

\begin{table}[h]\centering\small
\begin{tabular}{lcccc}
\toprule
 & neutral & glyph & dot & headings \\
\midrule
Untrained base model & .058 & .397 & \textbf{.434} & .338 \\
\bottomrule
\end{tabular}
\caption{Untrained Qwen2.5-7B-Instruct under the frozen protocol ($n=7{,}319$ per condition).}
\label{tab:base}
\end{table}

\subsection{Perturbations: identity matters, position does not}

Perturbing the glyph-trained model's prompt (Table~\ref{tab:pert}) reveals a sharp asymmetry: removing or substituting the glyphs collapses accuracy, while scrambling their positions or order costs nothing.
A model using markers as \emph{phase} structure should care about position; a model using them as a \emph{trigger} should care only about presence.
The data say trigger.

\begin{table}[h]\centering\small
\begin{tabular}{lcc}
\toprule
Perturbation & Accuracy & $\Delta$ vs.\ matched (.294) \\
\midrule
delete all glyphs        & .148 & $-14.6$ \\
replace with dots        & .176 & $-11.8$ \\
replace with unseen symbols & .191 & $-10.3$ \\
\midrule
cluster at prompt start  & .309 & $+1.5$ \\
displace into words      & .322 & $+2.7$ \\
permute order            & .308 & $+1.4$ \\
\bottomrule
\end{tabular}
\caption{Glyph-arm accuracy under prompt perturbations ($n=7{,}319$ each).}
\label{tab:pert}
\end{table}

\section{Mechanism: a rare-token mode key}

\paragraph{Tokenization.}
Each glyph encodes as three byte-fallback tokens sharing a two-token prefix; the four glyphs differ only in their final byte.
Every harmful perturbation removes these rare bytes; every harmless one preserves the exact bag of tokens.

\paragraph{Dose--response.}
Sufficiency tests on the glyph-trained model ($n=7{,}319$ each): four copies of a \emph{single} trained glyph nearly restore matched performance (27.2\% vs.\ 29.7\%); one lone glyph recovers most of the collapse (24.0\%); an untrained alchemical glyph sharing only the byte prefix gives partial credit (21.5\%); nothing gives 14.8\%.
The distinct four-glyph ``phase vocabulary'' is worth at most 2.5 points.
The trigger is graded rare-byte familiarity.

\paragraph{Weights.}
All six adapters have nearly equal update magnitude, but the \arm{glyph} and \arm{random} adapters share almost twice the subspace (mean principal-angle cosine 0.40) of any other pair (0.26--0.29): the two arms exposed to alchemical byte tokens learned the same distinct machinery, concentrated in late layers (peak 23--26 of 28).
Notably, \arm{glyph} reached the \emph{lowest} training loss of the reasoning arms while generalizing worst---the easy-to-predict marker tokens flatter the loss while the model degrades on the mathematics.

\paragraph{The key is a steerable direction.}
Let $d_\ell$ be the mean difference in the layer-$\ell$ residual stream at the final prompt token between glyph-present and glyph-deleted prompts.
Adding $d_{20}$ during generation on glyph-\emph{free} prompts lifts the glyph-trained model from 26.7\% to 38.7\% ($n=300$; matched-prompt reference 42.0\%).
A layer sweep ($n=100$) shows the effect is flat through layer~12, rises from layer~14, and peaks at layers 22--26, where steering reaches \textbf{56\%---fourteen points above the matched-glyph prompt itself}.
Masking the glyph positions from all attention \emph{after} prompt processing costs nothing (41\% vs.\ 42\%): the key is absorbed into the residual stream while the prompt is read and never re-read from the glyph tokens.

\paragraph{Why the prompt benefit did not survive training.}
The trained direction does nothing for the untrained base model (8.3\%, at baseline), and even the base model's own glyph-vs-neutral direction captures only a small fraction of its prompt effect ($+3.7$ of $\sim{+}30$ points).
Prompting therefore works through rich, distributed in-context processing that no single direction summarizes, whereas fine-tuning \emph{compressed} that behavior into a low-dimensional switch keyed to rare tokens.
The compression is where the capability went.

\section{Conclusion}

Under a controlled protocol, symbolic phase markers do not improve mathematical reasoning.
Structured prompts genuinely elicit better-formatted, more systematic answers from the untrained model, but arbitrary dots do this at least as well as glyphs.
Fine-tuning on glyph-formatted traces converts that free elicitation benefit into a brittle dependency: a single late-layer activation direction, triggered by rare byte tokens, that gates the trained behavior while overall capability degrades.
The constructive corollary is that the direction itself is a better handle than the symbols ever were---injecting it directly outperforms glyph prompting by fourteen points.

\paragraph{Limitations.}
All results are from the excluded pilot seed (42); confirmatory seeds have not been run.
The preregistered statistics stage (with its 300-item manual audit of the automatic parser) has not been executed, so no confidence intervals are reported.
The mechanism experiments (steering, layer sweep, KV masking, weight analysis) are exploratory, run on 100--300-problem subsets, and were not preregistered.
The untrained model's neutral-prompt score conflates reasoning failure with format noncompliance.

\end{document}
